\documentclass[12pt,a4paper]{article}

% --- Packages ---
\usepackage[utf8]{inputenc}
\usepackage[T1]{fontenc}
\usepackage[french]{babel}
\usepackage{geometry}
\usepackage{enumitem}
\usepackage{array}
\usepackage{fancyhdr}
\usepackage{amsmath, amssymb}

% --- Mise en page ---
\geometry{margin=2.5cm}

% --- Pied de page ---
\pagestyle{fancy}
\fancyhf{}
\fancyfoot[R]{\it{G.M.}}
\renewcommand{\headrulewidth}{0pt}
\renewcommand{\footrulewidth}{0pt}

\setlength{\parindent}{0pt}

\begin{document}

\begin{center}
    \Large \textbf{-- Automatismes - QCM--} \\[5pt]
    \normalsize Aucune justification n’est demandée et une seule réponse est possible par question.
\end{center}

\vspace{0.7cm}

\textbf{Question 1 :}
\vspace{5pt}

Dans une classe de 27 élèves, 9 sont gauchers. La proportion d’élèves gauchers est :
\vspace{0.3cm}

\renewcommand{\arraystretch}{3}
\begin{tabular}{|m{3.4cm}|m{3.4cm}|m{3.4cm}|m{3.4cm}|}
\hline
\centering \vspace{-0.3cm} \textbf{A.} 0,9 & 
\centering \vspace{-0.3cm} \textbf{B.} 0,3 & 
\centering \vspace{-0.3cm} \textbf{C.} 0,25 & 
\centering \vspace{-0.3cm} \textbf{D.} $\dfrac{1}{3}$ \tabularnewline
\hline
\end{tabular}

\vspace{0.7cm}

\textbf{Question 2 :}
\vspace{5pt}

Un article coûte 200 €. Le vendeur propose une réduction de 15\%. Combien d'euros représente la réduction ?
\vspace{0.3cm}

\renewcommand{\arraystretch}{3}
\begin{tabular}{|m{3.4cm}|m{3.4cm}|m{3.4cm}|m{3.4cm}|}
\hline
\centering \vspace{-0.3cm} \textbf{A.} 20 € & 
\centering \vspace{-0.3cm} \textbf{B.} 30 € & 
\centering \vspace{-0.3cm} \textbf{C.} 25 € & 
\centering \vspace{-0.3cm} \textbf{D.} 15 € \tabularnewline
\hline
\end{tabular}

\vspace{0.7cm}

\textbf{Question 3 :}
\vspace{5pt}

Dans un collège, 20\% des élèves sont en classe de 3ème, ce qui représente 50 élèves. L’effectif total du collège est :
\vspace{0.3cm}

\renewcommand{\arraystretch}{3}
\begin{tabular}{|m{3.4cm}|m{3.4cm}|m{3.4cm}|m{3.4cm}|}
\hline
\centering \vspace{-0.3cm} \textbf{A.} 250 & 
\centering \vspace{-0.3cm} \textbf{B.} 200 & 
\centering \vspace{-0.3cm} \textbf{C.} 300 & 
\centering \vspace{-0.3cm} \textbf{D.} 350 \tabularnewline
\hline
\end{tabular}

\vspace{0.7cm}

\textbf{Question 4 :}
\vspace{5pt}

Un prix augmente de 20\%. Le coefficient multiplicateur associé à cette augmentation est :

\vspace{0.3cm}

\renewcommand{\arraystretch}{3}
\begin{tabular}{|m{3.4cm}|m{3.4cm}|m{3.4cm}|m{3.4cm}|}
\hline
\centering \vspace{-0.3cm} \textbf{A.} 0,2 & 
\centering \vspace{-0.3cm} \textbf{B.} 0,8 & 
\centering \vspace{-0.3cm} \textbf{C.} 1,2 & 
\centering \vspace{-0.3cm} \textbf{D.} 1,25 \tabularnewline
\hline
\end{tabular}

\vspace{0.7cm}

\textbf{Question 5 :}
\vspace{5pt}

Un article coûte 120 €. Après une augmentation de 15\%, son nouveau prix est :
\vspace{0.3cm}

\renewcommand{\arraystretch}{3}
\begin{tabular}{|m{3.4cm}|m{3.4cm}|m{3.4cm}|m{3.4cm}|}
\hline
\centering \vspace{-0.3cm} \textbf{A.} 140 € & 
\centering \vspace{-0.3cm} \textbf{B.} 128 € & 
\centering \vspace{-0.3cm} \textbf{C.} 138 € & 
\centering \vspace{-0.3cm} \textbf{D.} 130 € \tabularnewline
\hline
\end{tabular}

\vspace{0.7cm}

\textbf{Question 6 :}
\vspace{5pt}

Un article coûte 180 €. Après une réduction de 20\%, son nouveau prix est :
\vspace{0.3cm}

\renewcommand{\arraystretch}{3}
\begin{tabular}{|m{3.4cm}|m{3.4cm}|m{3.4cm}|m{3.4cm}|}
\hline
\centering \vspace{-0.3cm} \textbf{A.} 144 € & 
\centering \vspace{-0.3cm} \textbf{B.} 150 € & 
\centering \vspace{-0.3cm} \textbf{C.} 140 € & 
\centering \vspace{-0.3cm} \textbf{D.} 136 € \tabularnewline
\hline
\end{tabular}

\vspace{0.7cm}

\textbf{Question 7 :}
\vspace{5pt}

Après une augmentation de 25\%, un article coûte 125 €. Son prix initial était :
\vspace{0.3cm}

\renewcommand{\arraystretch}{3}
\begin{tabular}{|m{3.4cm}|m{3.4cm}|m{3.4cm}|m{3.4cm}|}
\hline
\centering \vspace{-0.3cm} \textbf{A.} 105 € & 
\centering \vspace{-0.3cm} \textbf{B.} 95 € & 
\centering \vspace{-0.3cm} \textbf{C.} 110 € & 
\centering \vspace{-0.3cm} \textbf{D.} 100 € \tabularnewline
\hline
\end{tabular}

\vspace{0.7cm}

\textbf{Question 8 :}
\vspace{5pt}

Après une réduction de 10\%, un article coûte 81 €. Son prix initial était :
\vspace{0.3cm}

\renewcommand{\arraystretch}{3}
\begin{tabular}{|m{3.4cm}|m{3.4cm}|m{3.4cm}|m{3.4cm}|}
\hline
\centering \vspace{-0.3cm} \textbf{A.} 100 € & 
\centering \vspace{-0.3cm} \textbf{B.} 95 € & 
\centering \vspace{-0.3cm} \textbf{C.} 90 € & 
\centering \vspace{-0.3cm} \textbf{D.} 105 € \tabularnewline
\hline
\end{tabular}

\vspace{0.7cm}

\textbf{Question 9 :}
\vspace{5pt}

Un prix passe de 60 € à 78 €. Le taux d’évolution, exprimé en pourcentage, est :
\vspace{0.3cm}

\renewcommand{\arraystretch}{3}
\begin{tabular}{|m{3.4cm}|m{3.4cm}|m{3.4cm}|m{3.4cm}|}
\hline
\centering \vspace{-0.3cm} \textbf{A.} 20\% & 
\centering \vspace{-0.3cm} \textbf{B.} 30\% & 
\centering \vspace{-0.3cm} \textbf{C.} 25\% & 
\centering \vspace{-0.3cm} \textbf{D.} 15\% \tabularnewline
\hline
\end{tabular}

\vspace{0.7cm}

\textbf{Question 10 :}
\vspace{5pt}

Un prix subit deux augmentations successives de 5\% puis de 15\%. Le taux d’évolution global, exprimé en pourcentage, est :
\vspace{0.3cm}

\renewcommand{\arraystretch}{3}
\begin{tabular}{|m{3.4cm}|m{3.4cm}|m{3.4cm}|m{3.4cm}|}
\hline
\centering \vspace{-0.3cm} \textbf{A.} 20\% & 
\centering \vspace{-0.3cm} \textbf{B.} 21,75\% & 
\centering \vspace{-0.3cm} \textbf{C.} 20,75\% & 
\centering \vspace{-0.3cm} \textbf{D.} 25\% \tabularnewline
\hline
\end{tabular}

\end{document}
